\section{Objetivos}

\subsection{Objetivo general}

Desarrollar y validar un analizador léxico para un subconjunto definido de Prolog, relacionando expresiones regulares y autómatas finitos con una implementación en Python capaz de generar tokens posicionados y diagnósticos recuperables.

\subsection{Objetivos específicos}

Los objetivos específicos consisten en definir las categorías admitidas mediante patrones, ejemplos, atributos y reglas de prioridad; modelar categorías representativas con autómatas finitos y documentar una conversión de AFN a AFD junto con su minimización; implementar el recorrido de la fuente, la tabla de lexemas y los diagnósticos de error; y comprobar el comportamiento del analizador mediante pruebas automatizadas y dos archivos completos de entrada.

\section{Alcance del analizador}

La especificación restringe los identificadores y dígitos a caracteres ASCII. Esta decisión mantiene explícitas las clases de transición del autómata y acota la implementación al subconjunto pedido, aunque SWI-Prolog admite identificadores Unicode más amplios \cite{swisyntax}. Los literales citados sí pueden contener otros caracteres, siempre que no introduzcan un salto de línea ni una secuencia de escape fuera de la convención adoptada.

El analizador produce tokens y errores léxicos, pero no determina si esos tokens forman una cláusula correcta. Tampoco interpreta precedencias, realiza unificación, administra ámbitos ni asigna significado a los literales. La Tabla~\ref{tab:alcance} resume esa frontera para evitar que los resultados se presenten como una implementación completa de Prolog.

\begin{table}[H]
\centering
\small
\caption{Decisiones de alcance del analizador.}
\label{tab:alcance}
\begin{tabularx}{\textwidth}{L{2.85cm}YY}
\toprule
\tableheader{Aspecto} & \tableheader{Incluido} & \tableheader{Excluido} \\
\midrule
Identificadores & Átomos iniciados en minúscula; variables iniciadas en mayúscula o guion bajo; \texttt{\_} como anónima. & Unicode fuera de literales citados. \\
\rowcolor{palegray} Números & Enteros, decimales con dígitos a ambos lados del punto y exponentes opcionales. & Bases no decimales, racionales, infinidades y NaN. \\
Texto & Átomos con comillas simples y cadenas dobles; escapes \texttt{\textbackslash n}, \texttt{\textbackslash t}, \texttt{\textbackslash r}, \texttt{\textbackslash\textbackslash}, \texttt{\textbackslash'} y \texttt{\textbackslash"}. & Saltos de línea sin escapar y otros escapes. \\
\rowcolor{palegray} Operadores & Cláusulas, consultas, unificación, comparación, aritmética y control señalados en la tarea. & Operadores definidos por el usuario. \\
Comentarios & \texttt{\%} hasta fin de línea y \texttt{/* ... */}. & Comentarios de bloque anidados. \\
\rowcolor{palegray} Procesamiento & Tokenización, posiciones, tabla de lexemas y recuperación léxica. & Sintaxis, precedencia, unificación, ámbitos y semántica. \\
\bottomrule
\end{tabularx}
\end{table}
